% claim.tex -- claim of record (assembled file; regenerate, do not edit).
% Assembly: claim-skeleton.tex + the shipped proof files, spliced by
% build-claim.py (deliverable-directory tooling, not shipped). The
% machine-readable theorem anchors are the "% LEAN:" lines. All Lean listings
% are byte-identical splices of the shipped proof files (comment-cleaned
% editions; code regions byte-identical to the frozen development of record --
% correspondence record: audit/comment-cleanup-diff.md in this deposit):
%   lst:stmt  <- the lean fence of proofs/nfact-primorial-square.statement.txt
%   lst:chain <- proofs/final-nfact-primorial-square.lean, lines 96-168
%   lst:tail  <- proofs/final-nfact-primorial-square.lean, lines 735-767
%   lst:main  <- proofs/final-nfact-primorial-square.lean, lines 771-1094
% Byte-exactness record: audit/listing-verify.txt (ships in this deposit).
\documentclass[11pt]{article}

\usepackage[T1]{fontenc}
\usepackage[utf8]{inputenc}
\usepackage{amsmath,amssymb,amsthm}
\usepackage{graphicx}
\usepackage{hyperref}
\usepackage{xurl}
\setlength{\emergencystretch}{3em}
% lstlean.tex — Lean style for the listings package.
% Lineage: Jeremy Avigad (2015), adapted from Assia Mahboubi's SSR style;
% inspired by Olivier Verdier's unixode.sty for the Unicode literate table.
% No CTAN package or upstream repo exists; papers carry this file in their
% source package (cap set arXiv:1907.01449, sphere eversion arXiv:2210.07746).
% Adaptation note: sorry/admit/sorryAx are highlighted in a dedicated error
% tier so that incomplete proofs stand out visually in listings.
\usepackage{listings}
\usepackage{xcolor}

\definecolor{keywordcolor}{rgb}{0.7, 0.1, 0.1}   % red keywords
\definecolor{tacticcolor}{rgb}{0.0, 0.1, 0.6}    % blue tactics
\definecolor{commentcolor}{rgb}{0.4, 0.4, 0.4}   % grey comments
\definecolor{symbolcolor}{rgb}{0.0, 0.1, 0.6}    % blue symbols
\definecolor{sortcolor}{rgb}{0.1, 0.5, 0.1}      % green sorts
\definecolor{errorcolor}{rgb}{0.6, 0, 0}         % dark red: sorry/admit
\definecolor{stringcolor}{rgb}{0.5, 0.1, 0.5}    % purple strings

\lstdefinelanguage{lean}{
  % Anything between $ becomes LaTeX math mode
  mathescape=false,
  % Comments may or not include Latex commands
  texcl=false,
  % keywords, list taken from lean-syntax.txt
  morekeywords=[1]{import, prelude, protected, private, noncomputable,
    definition, meta, renaming, hiding, parameter, parameters, begin,
    constant, constants, lemma, variable, variables, theory, print,
    theorem, example, open, section, namespace, universe, universes, end,
    modifier, modifiers, using, export, exposing, axiom, inductive,
    coinductive, with, structure, class, instance, attribute, local, where,
    by, have, show, let, in, if, then, else, match, do, for, fun, assume,
    from, take, calc, include, omit, infix, infixl, infixr, notation,
    macro, syntax, elab, set_option, initialize, deriving, opaque, abbrev,
    mutual, partial, scoped, termination_by, decreasing_by},
  % Sorts
  morekeywords=[2]{Sort, Type, Prop},
  % tactics, list taken from lean-syntax.txt (tactic keywords)
  morekeywords=[3]{intro, intros, apply, exact, refine, rfl, simp, simp_all,
    simpa, simp_rw, rw, rewrite, erewrite, ring, ring_nf, omega, decide,
    norm_num, norm_cast, induction, cases, rcases, obtain, constructor,
    ext, funext, conv, congr, field_simp, linarith, nlinarith, positivity,
    tauto, trivial, aesop, use, by_cases, by_contra, push_neg, contrapose,
    symm, trans, assumption, contradiction, exfalso, specialize,
    generalize, revert, clear, rename_i, subst, unfold, delta, change,
    convert, split_ifs, interval_cases, fin_cases, zify, gcongr,
    nth_rewrite, suffices, generalizing, at, only, native_decide},
  % warnings - constitution tier: must never survive into a final proof
  morekeywords=[4]{sorry, admit, sorryAx},
  literate=
    {α}{{\ensuremath{\alpha}}}1
    {β}{{\ensuremath{\beta}}}1
    {γ}{{\ensuremath{\gamma}}}1
    {δ}{{\ensuremath{\delta}}}1
    {ε}{{\ensuremath{\varepsilon}}}1
    {ζ}{{\ensuremath{\zeta}}}1
    {η}{{\ensuremath{\eta}}}1
    {θ}{{\ensuremath{\theta}}}1
    {ι}{{\ensuremath{\iota}}}1
    {κ}{{\ensuremath{\kappa}}}1
    {λ}{{\ensuremath{\lambda}}}1
    {μ}{{\ensuremath{\mu}}}1
    {ν}{{\ensuremath{\nu}}}1
    {ξ}{{\ensuremath{\xi}}}1
    {π}{{\ensuremath{\pi}}}1
    {ρ}{{\ensuremath{\rho}}}1
    {σ}{{\ensuremath{\sigma}}}1
    {τ}{{\ensuremath{\tau}}}1
    {υ}{{\ensuremath{\upsilon}}}1
    {φ}{{\ensuremath{\varphi}}}1
    {χ}{{\ensuremath{\chi}}}1
    {ψ}{{\ensuremath{\psi}}}1
    {ω}{{\ensuremath{\omega}}}1
    {Γ}{{\ensuremath{\Gamma}}}1
    {Δ}{{\ensuremath{\Delta}}}1
    {Θ}{{\ensuremath{\Theta}}}1
    {Λ}{{\ensuremath{\Lambda}}}1
    {Ξ}{{\ensuremath{\Xi}}}1
    {Π}{{\ensuremath{\Pi}}}1
    {Σ}{{\ensuremath{\Sigma}}}1
    {Φ}{{\ensuremath{\Phi}}}1
    {Ψ}{{\ensuremath{\Psi}}}1
    {Ω}{{\ensuremath{\Omega}}}1
    {ℵ}{{\ensuremath{\aleph}}}1
    {∀}{{\ensuremath{\forall}}}1
    {∃!}{{\ensuremath{\exists}!}}1
    {∃}{{\ensuremath{\exists}}}1
    {∈}{{\ensuremath{\in}}}1
    {∉}{{\ensuremath{\notin}}}1
    {∘}{{\ensuremath{\circ}}}1
    {∩}{{\ensuremath{\cap}}}1
    {∪}{{\ensuremath{\cup}}}1
    {⊆}{{\ensuremath{\subseteq}}}1
    {⊂}{{\ensuremath{\subset}}}1
    {⊇}{{\ensuremath{\supseteq}}}1
    {⊃}{{\ensuremath{\supset}}}1
    {⊄}{{\ensuremath{\nsubseteq}}}1
    {∅}{{\ensuremath{\emptyset}}}1
    {∧}{{\ensuremath{\wedge}}}1
    {∨}{{\ensuremath{\vee}}}1
    {¬}{{\ensuremath{\neg}}}1
    {⊤}{{\ensuremath{\top}}}1
    {⊥}{{\ensuremath{\bot}}}1
    {↔}{{\ensuremath{\leftrightarrow}}}1
    {→}{{\ensuremath{\rightarrow}}}1
    {←}{{\ensuremath{\leftarrow}}}1
    {↑}{{\ensuremath{\uparrow}}}1
    {⇒}{{\ensuremath{\Rightarrow}}}1
    {⇐}{{\ensuremath{\Leftarrow}}}1
    {⟹}{{\ensuremath{\Longrightarrow}}}1
    {↦}{{\ensuremath{\mapsto}}}1
    {≤}{{\ensuremath{\leq}}}1
    {≥}{{\ensuremath{\geq}}}1
    {≠}{{\ensuremath{\neq}}}1
    {≈}{{\ensuremath{\approx}}}1
    {≡}{{\ensuremath{\equiv}}}1
    {≃}{{\ensuremath{\simeq}}}1
    {≅}{{\ensuremath{\cong}}}1
    {×}{{\ensuremath{\times}}}1
    {·}{{\ensuremath{\cdot}}}1
    {±}{{\ensuremath{\pm}}}1
    {⊕}{{\ensuremath{\oplus}}}1
    {⊗}{{\ensuremath{\otimes}}}1
    {⋆}{{\ensuremath{\star}}}1
    {▸}{{\ensuremath{\blacktriangleright}}}1
    {⟨}{{\ensuremath{\langle}}}1
    {⟩}{{\ensuremath{\rangle}}}1
    {⦃}{{\ensuremath{\{\!\mid}}}1
    {⦄}{{\ensuremath{\mid\!\}}}}1
    {⟦}{{\ensuremath{[\![}}}1
    {⟧}{{\ensuremath{]\!]}}}1
    {⌊}{{\ensuremath{\lfloor}}}1
    {⌋}{{\ensuremath{\rfloor}}}1
    {⌈}{{\ensuremath{\lceil}}}1
    {⌉}{{\ensuremath{\rceil}}}1
    {√}{{\ensuremath{\surd}}}1
    {∣}{{\ensuremath{\mid}}}1
    {∤}{{\ensuremath{\nmid}}}1
    {∎}{{\ensuremath{\blacksquare}}}1
    {⋯}{{\ensuremath{\cdots}}}1
    {…}{{\ensuremath{\ldots}}}1
    {∑}{{\ensuremath{\sum}}}1
    {∏}{{\ensuremath{\prod}}}1
    {⁻¹}{{\ensuremath{^{-1}}}}1
    {¹}{{\ensuremath{^1}}}1
    {²}{{\ensuremath{^2}}}1
    {³}{{\ensuremath{^3}}}1
    {ⁿ}{{\ensuremath{^n}}}1
    {ᵐ}{{\ensuremath{^m}}}1
    {₀}{{\ensuremath{_0}}}1
    {₁}{{\ensuremath{_1}}}1
    {₂}{{\ensuremath{_2}}}1
    {₃}{{\ensuremath{_3}}}1
    {₄}{{\ensuremath{_4}}}1
    {₅}{{\ensuremath{_5}}}1
    {ₙ}{{\ensuremath{_n}}}1
    {ₘ}{{\ensuremath{_m}}}1
    {ₖ}{{\ensuremath{_k}}}1
    {ᵢ}{{\ensuremath{_i}}}1
    {ⱼ}{{\ensuremath{_j}}}1
    {ℤ}{{\ensuremath{\mathbb{Z}}}}1
    {ℕ}{{\ensuremath{\mathbb{N}}}}1
    {ℚ}{{\ensuremath{\mathbb{Q}}}}1
    {ℝ}{{\ensuremath{\mathbb{R}}}}1
    {ℂ}{{\ensuremath{\mathbb{C}}}}1
    {𝔽}{{\ensuremath{\mathbb{F}}}}1
    {𝒫}{{\ensuremath{\mathcal{P}}}}1
    {∗}{{\ensuremath{\ast}}}1
    {•}{{\ensuremath{\bullet}}}1,
  % Comments
  morecomment=[l]{--},
  morecomment=[s]{/-}{-/},
  morestring=[b]",
  stringstyle=\color{stringcolor},
  keepspaces=true,
  keywordstyle=[1]\color{keywordcolor}\bfseries,
  keywordstyle=[2]\color{sortcolor}\bfseries,
  keywordstyle=[3]\color{tacticcolor}\bfseries,
  keywordstyle=[4]\color{errorcolor}\bfseries\underline,
  escapeinside={(*@}{@*)},
  columns=fullflexible,
  basicstyle=\ttfamily\small,
  commentstyle=\color{commentcolor}\itshape,
  breaklines=true,
  showstringspaces=false,
  frame=single,
  framesep=2mm,
  xleftmargin=4mm,
  numbers=left,
  numberstyle=\tiny\color{commentcolor},
  numbersep=6pt,
}

% Inline Lean code in prose: \lean{Int.ModEq.pow}
\newcommand{\lean}[1]{\lstinline[language=lean,basicstyle=\ttfamily\normalsize]|#1|}

\lstset{language=lean}

% --- local rendering patch (2026-10-06) --- tectonic/XeTeX does not apply
% the listings literate table to multi-byte characters, so every non-ASCII,
% non-CJK codepoint that occurs inside this note's listings is mapped to a
% math glyph as an active character via the standard \lccode/\lowercase idiom
% (ASCII-only source, no extra packages). Mapped set: 00A7 00AC 00B2 03BD
% 2115 2192 2194 2200 2203 2208 2211 2212 2223 2224 2227 2260 2264 2265 2270
% 230A 230B 27E8 27E9 27F8 27F9 27FA. Deliberately NOT mapped: U+00B7 (the
% Lean case bullet, which occurs in the proof code) --- it routes through
% xeCJK + SimSun (making it active first aborts the run, observed in earlier
% claim notes). Rendering only: listing source bytes are untouched, so the
% byte-exactness checks are unaffected, and the \ifdefined guard leaves the
% pdflatex path unaffected.
\ifdefined\XeTeXversion
  {\lccode`\~="00A7 \lowercase{\global\catcode`~=\active \global\def~{\S}}}
  {\lccode`\~="00AC \lowercase{\global\catcode`~=\active \global\def~{\ensuremath{\neg}}}}
  {\lccode`\~="00B2 \lowercase{\global\catcode`~=\active \global\def~{\ensuremath{{}^2}}}}
  {\lccode`\~="03BD \lowercase{\global\catcode`~=\active \global\def~{\ensuremath{\nu}}}}
  {\lccode`\~="2115 \lowercase{\global\catcode`~=\active \global\def~{\ensuremath{\mathbb{N}}}}}
  {\lccode`\~="2192 \lowercase{\global\catcode`~=\active \global\def~{\ensuremath{\rightarrow}}}}
  {\lccode`\~="2194 \lowercase{\global\catcode`~=\active \global\def~{\ensuremath{\leftrightarrow}}}}
  {\lccode`\~="2200 \lowercase{\global\catcode`~=\active \global\def~{\ensuremath{\forall}}}}
  {\lccode`\~="2203 \lowercase{\global\catcode`~=\active \global\def~{\ensuremath{\exists}}}}
  {\lccode`\~="2208 \lowercase{\global\catcode`~=\active \global\def~{\ensuremath{\in}}}}
  {\lccode`\~="2211 \lowercase{\global\catcode`~=\active \global\def~{\ensuremath{\sum}}}}
  {\lccode`\~="2212 \lowercase{\global\catcode`~=\active \global\def~{\ensuremath{-}}}}
  {\lccode`\~="2223 \lowercase{\global\catcode`~=\active \global\def~{\ensuremath{\mid}}}}
  {\lccode`\~="2224 \lowercase{\global\catcode`~=\active \global\def~{\ensuremath{\nmid}}}}
  {\lccode`\~="2227 \lowercase{\global\catcode`~=\active \global\def~{\ensuremath{\wedge}}}}
  {\lccode`\~="2260 \lowercase{\global\catcode`~=\active \global\def~{\ensuremath{\neq}}}}
  {\lccode`\~="2264 \lowercase{\global\catcode`~=\active \global\def~{\ensuremath{\leq}}}}
  {\lccode`\~="2265 \lowercase{\global\catcode`~=\active \global\def~{\ensuremath{\geq}}}}
  {\lccode`\~="2270 \lowercase{\global\catcode`~=\active \global\def~{\ensuremath{\nleq}}}}
  {\lccode`\~="230A \lowercase{\global\catcode`~=\active \global\def~{\ensuremath{\lfloor}}}}
  {\lccode`\~="230B \lowercase{\global\catcode`~=\active \global\def~{\ensuremath{\rfloor}}}}
  {\lccode`\~="27E8 \lowercase{\global\catcode`~=\active \global\def~{\ensuremath{\langle}}}}
  {\lccode`\~="27E9 \lowercase{\global\catcode`~=\active \global\def~{\ensuremath{\rangle}}}}
  {\lccode`\~="27F8 \lowercase{\global\catcode`~=\active \global\def~{\ensuremath{\Longleftarrow}}}}
  {\lccode`\~="27F9 \lowercase{\global\catcode`~=\active \global\def~{\ensuremath{\Longrightarrow}}}}
  {\lccode`\~="27FA \lowercase{\global\catcode`~=\active \global\def~{\ensuremath{\Longleftrightarrow}}}}
\fi
\ifdefined\XeTeXversion
  \usepackage{xeCJK}
  \setCJKmainfont{SimSun}
\fi

\newtheorem{theorem}{Theorem}
\newtheorem{lemma}[theorem]{Lemma}

\title{The Diophantine equation $x^2 = n! \cdot n\#$: the complete solution set\\
{\large (Lean 4 machine-checked claim of record)}}
\author{Aurora0134\thanks{Contact: via the GitHub account Aurora0134.}}
\date{October 6, 2026}

\begin{document}
\maketitle

\begin{abstract}
We prove that the Diophantine equation $x^2 = n! \cdot n\#$ (the factorial
times the primorial) has, for $n \in \mathbb{N}$, exactly the solutions
$n \in \{0,1,2,3,4,5\}$, with $x = 1, 1, 2, 6, 12, 60$. This settles one of
the two cases that Novakovi\'c (2026) recorded as believed-finite with the
words ``at the moment we were not able to give an argument''
\cite{novakovic2026}. The result is formalized in Lean~4 against mathlib
v4.34.0: every frozen statement compiles with no \texttt{sorry}, and
\texttt{\#print axioms} reports only \{propext, \texttt{Quot.sound},
\texttt{Classical.choice}\}. A novelty review (full query log shipped in
this deposit) found no occupying record in the queried channels. This note is
a priority
claim of record, not a full paper; the Lean sources and the audit and novelty
evidence ship inside this deposit.
\end{abstract}

\section{Introduction}
This note records a complete-solution-set theorem formalized in Lean~4,
with the kernel as the sole adjudicator at every step: the statements were
frozen by hash before proving, and a separate final audit re-compared them
against the frozen snapshot, scanned the whole proof file for
\texttt{sorry}/\texttt{admit} including comments, and re-printed the axiom
dependencies \cite{lean4,mathlib}. The object is the Diophantine
equation of the title, one of the concretely open cases listed by Novakovi\'c
\cite{novakovic2026}. The note is a priority claim of record, deposited for
timestamping; a narrative paper is prepared separately and will reference
this deposit. No literature survey is included---that expectation belongs to
arXiv moderation and does not apply to a Zenodo deposit.

\section{Statement}\label{sec:statement}

We write $n!$ for the factorial and $n\# = \prod_{p \le n,\ p \text{ prime}} p$
for the primorial; $\nu_p(m)$ denotes the exponent of the prime $p$ in $m$.
In the formalization, \texttt{primorial} is the top-level mathlib declaration
and $\nu_p(m)$ is \texttt{Nat.factorization m p}. Division in
Lemma~\ref{lem:B} is natural-number division, so the interval reads
$\lfloor n/3 \rfloor < p \le \lfloor n/2 \rfloor$; for prime existence this
coincides with the real reading, since $p$ is an integer.

% LEAN: tasks/20261006-nfact-primorial-square :: isSquare_factorial_mul_primorial_iff grade=完全证明
\begin{theorem}[complete solution set, iff form]\label{thm:main}
For every $n \in \mathbb{N}$, $n! \cdot n\#$ is a perfect square if and only
if $n \le 5$.
\end{theorem}

% LEAN: tasks/20261006-nfact-primorial-square :: exists_sq_factorial_mul_primorial_iff grade=完全证明
\begin{theorem}[existence-witness form]\label{thm:exists}
For every $n \in \mathbb{N}$, the equation $x^2 = n! \cdot n\#$ has a solution
$x \in \mathbb{N}$ if and only if $n \in \{0,1,2,3,4,5\}$.
\end{theorem}

% LEAN: tasks/20261006-nfact-primorial-square :: isSquare_iff_all_factorization_odd grade=完全证明
\begin{lemma}[valuation criterion]\label{lem:A}
For $n \ge 1$: $n! \cdot n\#$ is a perfect square if and only if
$\nu_p(n!)$ is odd for every prime $p \le n$.
\end{lemma}

% LEAN: tasks/20261006-nfact-primorial-square :: exists_prime_between_div_three_and_div_two grade=完全证明
\begin{lemma}[interval prime lemma]\label{lem:B}
For every integer $n \ge 22$ there exists a prime $p$ with $n/3 < p \le n/2$.
\end{lemma}

% LEAN: tasks/20261006-nfact-primorial-square :: witnesses_zero_to_five grade=完全证明
\begin{theorem}[solution witnesses]\label{thm:witnesses}
$1^2 = 0! \cdot 0\#$, $1^2 = 1! \cdot 1\#$, $2^2 = 2! \cdot 2\#$,
$6^2 = 3! \cdot 3\#$, $12^2 = 4! \cdot 4\#$, and $60^2 = 5! \cdot 5\#$.
\end{theorem}

\noindent Three auxiliary factorization statements belong to the same frozen
set and carry the same grade (verbatim in Listing~\ref{lst:stmt}):
$\nu_p(n\#) = 1$ for primes $p \le n$, $\nu_p(n\#) = 0$ for primes
$p \nleq n$, and $\nu_p(n! \cdot n\#) = \nu_p(n!) + 1$ for primes $p \le n$.
% LEAN: tasks/20261006-nfact-primorial-square :: primorial_factorization_eq_one grade=完全证明
% LEAN: tasks/20261006-nfact-primorial-square :: primorial_factorization_eq_zero grade=完全证明
% LEAN: tasks/20261006-nfact-primorial-square :: factorization_factorial_mul_primorial grade=完全证明

\begin{lstlisting}[caption={Frozen statement snapshot (Lean 4): the complete
frozen set, copied from the \texttt{lean} fence of the snapshot shipped as
\texttt{proofs/nfact-primorial-square.statement.txt} (a comment-cleaned
edition: docstrings translated to English for publication; the statement
text is byte-identical to the frozen snapshot of record, sha256
\texttt{0b13d08d}; correspondence record
\texttt{audit/comment-cleanup-diff.md}).},label={lst:stmt}]
/-- **Main theorem (complete solution set, iff form)**:
    `n! * n#` is a perfect square ⟺ `n ≤ 5` (for `n : ℕ`).
    Over ℕ, `n ≤ 5 ⟺ n ∈ {0,1,2,3,4,5}`, so this is the main statement "the solutions are exactly n ∈ {0,...,5}". -/
theorem isSquare_factorial_mul_primorial_iff (n : ℕ) :
    IsSquare (n ! * primorial n) ↔ n ≤ 5

/-- **Main theorem (existence-witness form, delivered alongside the iff form)**:
    the Diophantine equation `x² = n! * n#` has a solution ⟺ `n ∈ {0,1,2,3,4,5}`. -/
theorem exists_sq_factorial_mul_primorial_iff (n : ℕ) :
    (∃ x : ℕ, x ^ 2 = n ! * primorial n) ↔ n ∈ ({0, 1, 2, 3, 4, 5} : Finset ℕ)

/-- **Lemma A (equivalent valuation criterion)**: for `n ≥ 1`, `n! * n#` is a perfect square ⟺
    `ν_p(n!)` is odd for every prime `p ≤ n`. -/
theorem isSquare_iff_all_factorization_odd (n : ℕ) (hn : 1 ≤ n) :
    IsSquare (n ! * primorial n) ↔
      ∀ p : ℕ, p.Prime → p ≤ n → Odd ((n !).factorization p)

/-- **Lemma B (interval prime lemma)**: for every integer `n ≥ 22` there exists a prime `p` with `n/3 < p ≤ n/2`. -/
theorem exists_prime_between_div_three_and_div_two (n : ℕ) (hn : 22 ≤ n) :
    ∃ p : ℕ, p.Prime ∧ n / 3 < p ∧ p ≤ n / 2

/-- **Witness layer (explicit square roots for the ⟸ direction of the main theorem)**: `x² = n! * n#` at each of n = 0..5, with x = 1, 1, 2, 6, 12, 60. -/
theorem witnesses_zero_to_five :
    1 ^ 2 = 0 ! * primorial 0 ∧
    1 ^ 2 = 1 ! * primorial 1 ∧
    2 ^ 2 = 2 ! * primorial 2 ∧
    6 ^ 2 = 3 ! * primorial 3 ∧
    12 ^ 2 = 4 ! * primorial 4 ∧
    60 ^ 2 = 5 ! * primorial 5

/-- Auxiliary statement: `ν_p(n#) = 1` (primes `p ≤ n`). -/
theorem primorial_factorization_eq_one {p n : ℕ} (hp : p.Prime) (hpn : p ≤ n) :
    (primorial n).factorization p = 1

/-- Auxiliary statement: `ν_p(n#) = 0` (primes `p ≰ n`). -/
theorem primorial_factorization_eq_zero {p n : ℕ} (hp : p.Prime) (hpn : ¬ p ≤ n) :
    (primorial n).factorization p = 0

/-- Auxiliary statement (the multiplicative skeleton of the criterion): `ν_p(n! * n#) = ν_p(n!) + 1` (primes `p ≤ n`). -/
theorem factorization_factorial_mul_primorial {p n : ℕ} (hp : p.Prime) (hpn : p ≤ n) :
    (n ! * primorial n).factorization p = (n !).factorization p + 1
\end{lstlisting}

\section{Proof}\label{sec:proof}

The proof assembles four components. (1) \emph{The valuation criterion}
(Lemma~\ref{lem:A}): $n\#$ contributes each prime $p \le n$ exactly once, so
$n! \cdot n\#$ is a square if and only if every $\nu_p(n!)$ is odd; both
directions go through the library characterization of squares by even
factorizations. (2) \emph{Exclusion of large $n$}: for $n \ge 22$,
Lemma~\ref{lem:B} yields a prime $p$ with $n/3 < p \le n/2$; then
$2p \le n < 3p$ and $p^2 > n$, so Legendre's formula gives
$\nu_p(n!) = \lfloor n/p \rfloor + \lfloor n/p^2 \rfloor + \cdots = 2$, an
even number, contradicting the all-odd criterion. Lemma~\ref{lem:B} itself is
proved in two regimes. For $22 \le n \le 999999$, a certificate chain of
twenty-nine primes (Listing~\ref{lst:chain}), in which each prime $q$ covers
$2q \le n \le 3q-1$ and consecutive certificates overlap, closes the range by
\texttt{norm\_num}/\texttt{omega}. For $n \ge 10^6$, an Erd\H{o}s-style
factorial-ratio argument (Listing~\ref{lst:tail}): with
$\mathrm{tailR}(n) = (n/2)!\,(n/6)!\,/\,((n/3)!\,(n/4)!\,(n/12)!)$ one has
$18^{\lfloor n/12 \rfloor} < \mathrm{tailR}(n) \cdot (4\lfloor n/12
\rfloor+3)^6$, while the absence of primes in $(n/3, n/2]$ would force
$\mathrm{tailR}(n)$ below essentially $16^{\lfloor n/12 \rfloor}$ times
lower-order factors; the decisive inequality $18 > 16$ is closed by a
concavity two-point estimate. The segment $10^6 \le n < 1.2 \cdot 10^6$ is
closed by a single certificate hop with $q = 400009$. (3) \emph{Small $n$}:
witnesses for $n = 0, \ldots, 5$ (Theorem~\ref{thm:witnesses}) and, for
$6 \le n \le 21$, sixteen explicit factorizations each exhibiting a prime
with odd exponent (Listing~\ref{lst:main}). (4) \emph{Assembly}: the two main
theorems follow from the criterion, the exclusion, and the witnesses.

For honesty of the record: the route originally envisaged for
Lemma~\ref{lem:B} (splicing mathlib's Bertrand theorem across two dyadic
sub-intervals) was strictly refuted during the task---a countermodel (the set
of all primes with $13$ removed) satisfies the Bertrand property yet leaves
the window $(11, 16]$ empty, so no black-box argument at ratio $2$ can close
a window of ratio $3/2$---and the shipped route is the certificate chain plus
factorial-ratio tail described above. The complete proof file (including the
twenty-nine prime certificates and the twenty-eight continuity certificates,
which are one-line \texttt{norm\_num}/\texttt{omega} facts) ships in
\texttt{proofs/}; the blocks below are verbatim excerpts, each a contiguous
line range of that file. The shipped file is the comment-cleaned edition; its
code is byte-identical to the frozen development (record:
\texttt{audit/comment-cleanup-diff.md}).

\begin{lstlisting}[basicstyle=\ttfamily\footnotesize,caption={Excerpt of the
proof file (lines 96--168 of \texttt{final-nfact-primorial-square.lean},
which ships in full in \texttt{proofs/}): the certificate-chain mechanism for
Lemma~\ref{lem:B} on $22 \le n \le 999999$ --- the two boundary certificates,
the single-hop lemma \texttt{lemmaB\_hop}, and the twenty-nine-hop decision
tree.},label={lst:chain}]
theorem chain_start : 22 = 2 * 11 := by omega
theorem chain_end : 999999 ≤ 3 * 409609 - 1 := by omega

/-- Certificate-chain single hop: for prime `q` and `n ∈ [2q, 3q - 1]`,
    `q` itself witnesses `n/3 < q ≤ n/2`.
    Note: extract `hq.two_le` first, otherwise `omega` gets stuck at the `q = 0` counterexample under saturating subtraction on ℕ. -/
theorem lemmaB_hop {q : ℕ} (hq : q.Prime) {n : ℕ} (h1 : 2 * q ≤ n) (h2 : n ≤ 3 * q - 1) :
    ∃ p : ℕ, p.Prime ∧ n / 3 < p ∧ p ≤ n / 2 :=
  have hq2 : 2 ≤ q := hq.two_le
  ⟨q, hq, by omega, by omega⟩

/-- **Lemma B, finite segment (certificate-chain decision tree)**: for every `n ∈ [22, 999999]`,
    `∃ p, p.Prime ∧ n/3 < p ∧ p ≤ n/2`.
    Joins the tail theorem (which starts at n = 1000000) by a case split on `n < 1000000`. -/
theorem lemmaB_upto_999999 (n : ℕ) (h1 : 22 ≤ n) (h2 : n ≤ 999999) :
    ∃ p : ℕ, p.Prime ∧ n / 3 < p ∧ p ≤ n / 2 := by
  rcases lt_or_ge n 33 with h | h
  · exact lemmaB_hop (q := 11) prime_cert_11 (by omega) (by omega)
  · rcases lt_or_ge n 39 with h | h
    · exact lemmaB_hop (q := 13) prime_cert_13 (by omega) (by omega)
    · rcases lt_or_ge n 57 with h | h
      · exact lemmaB_hop (q := 19) prime_cert_19 (by omega) (by omega)
      · rcases lt_or_ge n 69 with h | h
        · exact lemmaB_hop (q := 23) prime_cert_23 (by omega) (by omega)
        · rcases lt_or_ge n 93 with h | h
          · exact lemmaB_hop (q := 31) prime_cert_31 (by omega) (by omega)
          · rcases lt_or_ge n 129 with h | h
            · exact lemmaB_hop (q := 43) prime_cert_43 (by omega) (by omega)
            · rcases lt_or_ge n 183 with h | h
              · exact lemmaB_hop (q := 61) prime_cert_61 (by omega) (by omega)
              · rcases lt_or_ge n 267 with h | h
                · exact lemmaB_hop (q := 89) prime_cert_89 (by omega) (by omega)
                · rcases lt_or_ge n 393 with h | h
                  · exact lemmaB_hop (q := 131) prime_cert_131 (by omega) (by omega)
                  · rcases lt_or_ge n 579 with h | h
                    · exact lemmaB_hop (q := 193) prime_cert_193 (by omega) (by omega)
                    · rcases lt_or_ge n 849 with h | h
                      · exact lemmaB_hop (q := 283) prime_cert_283 (by omega) (by omega)
                      · rcases lt_or_ge n 1263 with h | h
                        · exact lemmaB_hop (q := 421) prime_cert_421 (by omega) (by omega)
                        · rcases lt_or_ge n 1893 with h | h
                          · exact lemmaB_hop (q := 631) prime_cert_631 (by omega) (by omega)
                          · rcases lt_or_ge n 2823 with h | h
                            · exact lemmaB_hop (q := 941) prime_cert_941 (by omega) (by omega)
                            · rcases lt_or_ge n 4227 with h | h
                              · exact lemmaB_hop (q := 1409) prime_cert_1409 (by omega) (by omega)
                              · rcases lt_or_ge n 6339 with h | h
                                · exact lemmaB_hop (q := 2113) prime_cert_2113 (by omega) (by omega)
                                · rcases lt_or_ge n 9507 with h | h
                                  · exact lemmaB_hop (q := 3169) prime_cert_3169 (by omega) (by omega)
                                  · rcases lt_or_ge n 14253 with h | h
                                    · exact lemmaB_hop (q := 4751) prime_cert_4751 (by omega) (by omega)
                                    · rcases lt_or_ge n 21363 with h | h
                                      · exact lemmaB_hop (q := 7121) prime_cert_7121 (by omega) (by omega)
                                      · rcases lt_or_ge n 32001 with h | h
                                        · exact lemmaB_hop (q := 10667) prime_cert_10667 (by omega) (by omega)
                                        · rcases lt_or_ge n 47973 with h | h
                                          · exact lemmaB_hop (q := 15991) prime_cert_15991 (by omega) (by omega)
                                          · rcases lt_or_ge n 71943 with h | h
                                            · exact lemmaB_hop (q := 23981) prime_cert_23981 (by omega) (by omega)
                                            · rcases lt_or_ge n 107907 with h | h
                                              · exact lemmaB_hop (q := 35969) prime_cert_35969 (by omega) (by omega)
                                              · rcases lt_or_ge n 161853 with h | h
                                                · exact lemmaB_hop (q := 53951) prime_cert_53951 (by omega) (by omega)
                                                · rcases lt_or_ge n 242769 with h | h
                                                  · exact lemmaB_hop (q := 80923) prime_cert_80923 (by omega) (by omega)
                                                  · rcases lt_or_ge n 364137 with h | h
                                                    · exact lemmaB_hop (q := 121379) prime_cert_121379 (by omega) (by omega)
                                                    · rcases lt_or_ge n 546177 with h | h
                                                      · exact lemmaB_hop (q := 182059) prime_cert_182059 (by omega) (by omega)
                                                      · rcases lt_or_ge n 819249 with h | h
                                                        · exact lemmaB_hop (q := 273083) prime_cert_273083 (by omega) (by omega)
                                                        · exact lemmaB_hop (q := 409609) prime_cert_409609 (by omega) (by omega)
\end{lstlisting}

\begin{lstlisting}[basicstyle=\ttfamily\footnotesize,caption={Excerpt of the
proof file (lines 735--767): the tail theorem for $n \ge 10^6$
(\texttt{lemmaB\_tail\_1000000}), closing Lemma~\ref{lem:B} on the infinite
regime via the factorial-ratio estimate.},label={lst:tail}]
theorem lemmaB_tail_1000000 (n : ℕ) (hn : 1000000 ≤ n) :
    ∃ p : ℕ, p.Prime ∧ n / 3 < p ∧ p ≤ n / 2 := by
  rcases lt_or_ge n 1200000 with h | h
  · exact lemmaB_hop (q := 400009) (by norm_num) (by omega) (by omega)
  · by_contra hc
    push_neg at hc
    have hc' : ∀ p : ℕ, p.Prime → n / 3 < p → ¬ p ≤ n / 2 := by
      intro p hp1 hp2 hle
      have hcc := hc p hp1 hp2
      omega
    have hub := tailR_le_of_no_prime n (by omega) hc'
    have hlb := tailR_lower_strict n hn
    have hmain := tail_main_inequality_nat (n/12) (by omega)
    have h2 : n < (n/12 + 1) * 12 := (Nat.div_lt_iff_lt_mul (by norm_num : (0:ℕ) < 12)).mp (Nat.lt_succ_self _)
    have e1 : (n/2)^(Nat.sqrt (n/2)) ≤ (6*(n/12)+5)^(Nat.sqrt (6*(n/12)+5)) := by
      have hbase : n/2 ≤ 6*(n/12)+5 := by omega
      have hs : Nat.sqrt (n/2) ≤ Nat.sqrt (6*(n/12)+5) := Nat.sqrt_le_sqrt hbase
      exact (Nat.pow_le_pow_right (by omega : 0 < n/2) hs).trans (Nat.pow_le_pow_left hbase _)
    have e2 : 4^(n/6) ≤ 4 * 16^(n/12) := by
      have h4 : n/6 ≤ 2*(n/12) + 1 := by omega
      calc 4^(n/6) ≤ 4^(2*(n/12)+1) := Nat.pow_le_pow_right (by norm_num) h4
        _ = 4 * 16^(n/12) := by
          rw [pow_succ, pow_mul]
          norm_num
          ring
    have key : 18^(n/12) < 18^(n/12) := by
      calc 18^(n/12) < tailR n * (4*(n/12)+3)^6 := hlb
        _ ≤ ((n/2)^(Nat.sqrt (n/2)) * 4^(n/6)) * (4*(n/12)+3)^6 := Nat.mul_le_mul_right _ hub
        _ ≤ ((6*(n/12)+5)^(Nat.sqrt (6*(n/12)+5)) * (4 * 16^(n/12))) * (4*(n/12)+3)^6 :=
            Nat.mul_le_mul_right _ (Nat.mul_le_mul e1 e2)
        _ = 4 * (4*(n/12)+3)^6 * (6*(n/12)+5)^(Nat.sqrt (6*(n/12)+5)) * 16^(n/12) := by ring
        _ ≤ 18^(n/12) := hmain
    exact absurd key (lt_irrefl _)
\end{lstlisting}

\begin{lstlisting}[basicstyle=\ttfamily\footnotesize,caption={Excerpt of the
proof file (lines 771--1094): the main namespace \texttt{NFactPrimorialSquare}
in full --- the three auxiliary factorization lemmas, Lemma~\ref{lem:A}, the
witness layer, the sixteen non-square witnesses, Lemma~\ref{lem:B} assembled
from the chain and the tail, and the two main theorems.},label={lst:main}]
namespace NFactPrimorialSquare

/-- Auxiliary statement: `ν_p(n#) = 1` for primes `p ≤ n`. -/
theorem primorial_factorization_eq_one {p n : ℕ} (hp : p.Prime) (hpn : p ≤ n) :
    (primorial n).factorization p = 1 :=
  factorization_eq_one_of_squarefree (squarefree_primorial n) hp
    ((Nat.Prime.dvd_primorial_iff hp).mpr hpn)

/-- Auxiliary statement: `ν_p(n#) = 0` for primes `p ≰ n`. -/
theorem primorial_factorization_eq_zero {p n : ℕ} (hp : p.Prime) (hpn : ¬ p ≤ n) :
    (primorial n).factorization p = 0 :=
  Nat.factorization_eq_zero_of_not_dvd
    (fun h => hpn ((Nat.Prime.dvd_primorial_iff hp).mp h))

/-- Auxiliary statement (the multiplicative skeleton of the criterion): `ν_p(n! * n#) = ν_p(n!) + 1` for primes `p ≤ n`. -/
theorem factorization_factorial_mul_primorial {p n : ℕ} (hp : p.Prime) (hpn : p ≤ n) :
    (n ! * primorial n).factorization p = (n !).factorization p + 1 := by
  rw [Nat.factorization_mul (Nat.factorial_ne_zero n) (primorial_ne_zero n),
      Finsupp.add_apply, primorial_factorization_eq_one hp hpn]

/-- **Lemma A (equivalent valuation criterion)**: for `n ≥ 1`, `n! * n#` is a perfect square ⟺
    `ν_p(n!)` is odd for every prime `p ≤ n`.
    (`Nat.factorization m p` is ν_p(m): the exponent of the prime `p` in `m`.) -/
theorem isSquare_iff_all_factorization_odd (n : ℕ) (hn : 1 ≤ n) :
    IsSquare (n ! * primorial n) ↔
      ∀ p : ℕ, p.Prime → p ≤ n → Odd ((n !).factorization p) := by
  constructor
  · intro h p hp hpn
    have hEven : Even ((n ! * primorial n).factorization p) :=
      Nat.isSquare_iff_even_factorization.mp h p hp
    rw [factorization_factorial_mul_primorial hp hpn] at hEven
    obtain ⟨m, hm⟩ := hEven
    cases m with
    | zero => simp at hm
    | succ k => exact ⟨k, by omega⟩
  · intro h
    apply Nat.isSquare_iff_even_factorization.mpr
    intro p hp
    by_cases hpn : p ≤ n
    · rw [factorization_factorial_mul_primorial hp hpn]
      exact (h p hp hpn).add_one
    · rw [Nat.factorization_mul (Nat.factorial_ne_zero n) (primorial_ne_zero n),
        Finsupp.add_apply,
        Nat.factorization_eq_zero_of_not_dvd
          (fun hd => hpn ((Nat.Prime.dvd_factorial hp).mp hd)),
        primorial_factorization_eq_zero hp hpn]
      exact ⟨0, rfl⟩

/-- **Witness layer (explicit square roots for the ⟸ direction of the main theorem)**: `x² = n! * n#`
    at each of n = 0..5, with x = 1, 1, 2, 6, 12, 60 (the pointwise correspondence of the main statement). -/
theorem witnesses_zero_to_five :
    1 ^ 2 = 0 ! * primorial 0 ∧
    1 ^ 2 = 1 ! * primorial 1 ∧
    2 ^ 2 = 2 ! * primorial 2 ∧
    6 ^ 2 = 3 ! * primorial 3 ∧
    12 ^ 2 = 4 ! * primorial 4 ∧
    60 ^ 2 = 5 ! * primorial 5 := by decide

/-- Bridge lemma (internal to the witness layer): an odd factorization at some prime implies non-square.
    `IsSquare` gives `m = r * r`, so `ν_p(m) = ν_p(r) + ν_p(r)` is even, contradicting the oddness assumption. -/
theorem not_isSquare_of_odd_factorization {m p e : ℕ} (hm : m ≠ 0)
    (h : m.factorization p = 2 * e + 1) : ¬ IsSquare m := by
  rintro ⟨r, hr⟩
  have hne : r ≠ 0 := by rintro rfl; rw [mul_zero] at hr; exact hm hr
  rw [hr, Nat.factorization_mul hne hne, Finsupp.add_apply] at h
  omega

/-- Non-square witness n = 6: `ν_2(6! * 6#) = 5` (odd); `6! * 6# = 2^5 * 675`, and `2 ∤ 675`. -/
theorem nonsquare_at_6 : ¬ IsSquare (6 ! * primorial 6) := by
  have hM : (6 : ℕ) ! * primorial 6 = 2 ^ 5 * 675 := by decide
  apply not_isSquare_of_odd_factorization (p := 2) (e := 2)
    (mul_ne_zero (Nat.factorial_ne_zero 6) (primorial_ne_zero 6))
  rw [hM, Nat.factorization_mul (by norm_num : ((2:ℕ)^5) ≠ 0) (by norm_num : (675:ℕ) ≠ 0),
    Finsupp.add_apply, Nat.factorization_eq_zero_of_not_dvd (by decide : ¬ (2:ℕ) ∣ 675),
    Nat.factorization_pow, Finsupp.smul_apply,
    Nat.Prime.factorization Nat.prime_two, Finsupp.single_eq_same]; norm_num [nsmul_eq_mul]

/-- Non-square witness n = 7: `ν_2(7! * 7#) = 5` (odd); `7! * 7# = 2^5 * 33075`, and `2 ∤ 33075`. -/
theorem nonsquare_at_7 : ¬ IsSquare (7 ! * primorial 7) := by
  have hM : (7 : ℕ) ! * primorial 7 = 2 ^ 5 * 33075 := by decide
  apply not_isSquare_of_odd_factorization (p := 2) (e := 2)
    (mul_ne_zero (Nat.factorial_ne_zero 7) (primorial_ne_zero 7))
  rw [hM, Nat.factorization_mul (by norm_num : ((2:ℕ)^5) ≠ 0) (by norm_num : (33075:ℕ) ≠ 0),
    Finsupp.add_apply, Nat.factorization_eq_zero_of_not_dvd (by decide : ¬ (2:ℕ) ∣ 33075),
    Nat.factorization_pow, Finsupp.smul_apply,
    Nat.Prime.factorization Nat.prime_two, Finsupp.single_eq_same]; norm_num [nsmul_eq_mul]

/-- Non-square witness n = 8: `ν_3(8! * 8#) = 3` (odd); `8! * 8# = 3^3 * 313600`, and `3 ∤ 313600`. -/
theorem nonsquare_at_8 : ¬ IsSquare (8 ! * primorial 8) := by
  have hM : (8 : ℕ) ! * primorial 8 = 3 ^ 3 * 313600 := by decide
  apply not_isSquare_of_odd_factorization (p := 3) (e := 1)
    (mul_ne_zero (Nat.factorial_ne_zero 8) (primorial_ne_zero 8))
  rw [hM, Nat.factorization_mul (by norm_num : ((3:ℕ)^3) ≠ 0) (by norm_num : (313600:ℕ) ≠ 0),
    Finsupp.add_apply, Nat.factorization_eq_zero_of_not_dvd (by decide : ¬ (3:ℕ) ∣ 313600),
    Nat.factorization_pow, Finsupp.smul_apply,
    Nat.Prime.factorization Nat.prime_three, Finsupp.single_eq_same]; norm_num [nsmul_eq_mul]

/-- Non-square witness n = 9: `ν_3(9! * 9#) = 5` (odd); `9! * 9# = 3^5 * 313600`, and `3 ∤ 313600`. -/
theorem nonsquare_at_9 : ¬ IsSquare (9 ! * primorial 9) := by
  have hM : (9 : ℕ) ! * primorial 9 = 3 ^ 5 * 313600 := by decide
  apply not_isSquare_of_odd_factorization (p := 3) (e := 2)
    (mul_ne_zero (Nat.factorial_ne_zero 9) (primorial_ne_zero 9))
  rw [hM, Nat.factorization_mul (by norm_num : ((3:ℕ)^5) ≠ 0) (by norm_num : (313600:ℕ) ≠ 0),
    Finsupp.add_apply, Nat.factorization_eq_zero_of_not_dvd (by decide : ¬ (3:ℕ) ∣ 313600),
    Nat.factorization_pow, Finsupp.smul_apply,
    Nat.Prime.factorization Nat.prime_three, Finsupp.single_eq_same]; norm_num [nsmul_eq_mul]

/-- Non-square witness n = 10: `ν_2(10! * 10#) = 9` (odd); `10! * 10# = 2^9 * 1488375`, and `2 ∤ 1488375`. -/
theorem nonsquare_at_10 : ¬ IsSquare (10 ! * primorial 10) := by
  have hM : (10 : ℕ) ! * primorial 10 = 2 ^ 9 * 1488375 := by decide
  apply not_isSquare_of_odd_factorization (p := 2) (e := 4)
    (mul_ne_zero (Nat.factorial_ne_zero 10) (primorial_ne_zero 10))
  rw [hM, Nat.factorization_mul (by norm_num : ((2:ℕ)^9) ≠ 0) (by norm_num : (1488375:ℕ) ≠ 0),
    Finsupp.add_apply, Nat.factorization_eq_zero_of_not_dvd (by decide : ¬ (2:ℕ) ∣ 1488375),
    Nat.factorization_pow, Finsupp.smul_apply,
    Nat.Prime.factorization Nat.prime_two, Finsupp.single_eq_same]; norm_num [nsmul_eq_mul]

/-- Non-square witness n = 11: `ν_2(11! * 11#) = 9` (odd); `11! * 11# = 2^9 * 180093375`, and `2 ∤ 180093375`. -/
theorem nonsquare_at_11 : ¬ IsSquare (11 ! * primorial 11) := by
  have hM : (11 : ℕ) ! * primorial 11 = 2 ^ 9 * 180093375 := by decide
  apply not_isSquare_of_odd_factorization (p := 2) (e := 4)
    (mul_ne_zero (Nat.factorial_ne_zero 11) (primorial_ne_zero 11))
  rw [hM, Nat.factorization_mul (by norm_num : ((2:ℕ)^9) ≠ 0) (by norm_num : (180093375:ℕ) ≠ 0),
    Finsupp.add_apply, Nat.factorization_eq_zero_of_not_dvd (by decide : ¬ (2:ℕ) ∣ 180093375),
    Nat.factorization_pow, Finsupp.smul_apply,
    Nat.Prime.factorization Nat.prime_two, Finsupp.single_eq_same]; norm_num [nsmul_eq_mul]

/-- Non-square witness n = 12: `ν_2(12! * 12#) = 11` (odd); `12! * 12# = 2^11 * 540280125`, and `2 ∤ 540280125`. -/
theorem nonsquare_at_12 : ¬ IsSquare (12 ! * primorial 12) := by
  have hM : (12 : ℕ) ! * primorial 12 = 2 ^ 11 * 540280125 := by decide
  apply not_isSquare_of_odd_factorization (p := 2) (e := 5)
    (mul_ne_zero (Nat.factorial_ne_zero 12) (primorial_ne_zero 12))
  rw [hM, Nat.factorization_mul (by norm_num : ((2:ℕ)^11) ≠ 0) (by norm_num : (540280125:ℕ) ≠ 0),
    Finsupp.add_apply, Nat.factorization_eq_zero_of_not_dvd (by decide : ¬ (2:ℕ) ∣ 540280125),
    Nat.factorization_pow, Finsupp.smul_apply,
    Nat.Prime.factorization Nat.prime_two, Finsupp.single_eq_same]; norm_num [nsmul_eq_mul]

/-- Non-square witness n = 13: `ν_2(13! * 13#) = 11` (odd); `13! * 13# = 2^11 * 91307341125`, and `2 ∤ 91307341125`. -/
theorem nonsquare_at_13 : ¬ IsSquare (13 ! * primorial 13) := by
  have hM : (13 : ℕ) ! * primorial 13 = 2 ^ 11 * 91307341125 := by decide
  apply not_isSquare_of_odd_factorization (p := 2) (e := 5)
    (mul_ne_zero (Nat.factorial_ne_zero 13) (primorial_ne_zero 13))
  rw [hM, Nat.factorization_mul (by norm_num : ((2:ℕ)^11) ≠ 0) (by norm_num : (91307341125:ℕ) ≠ 0),
    Finsupp.add_apply, Nat.factorization_eq_zero_of_not_dvd (by decide : ¬ (2:ℕ) ∣ 91307341125),
    Nat.factorization_pow, Finsupp.smul_apply,
    Nat.Prime.factorization Nat.prime_two, Finsupp.single_eq_same]; norm_num [nsmul_eq_mul]

/-- Non-square witness n = 14: `ν_5(14! * 14#) = 3` (odd); `14! * 14# = 5^3 * 20943712677888`, and `5 ∤ 20943712677888`. -/
theorem nonsquare_at_14 : ¬ IsSquare (14 ! * primorial 14) := by
  have hM : (14 : ℕ) ! * primorial 14 = 5 ^ 3 * 20943712677888 := by decide
  apply not_isSquare_of_odd_factorization (p := 5) (e := 1)
    (mul_ne_zero (Nat.factorial_ne_zero 14) (primorial_ne_zero 14))
  rw [hM, Nat.factorization_mul (by norm_num : ((5:ℕ)^3) ≠ 0) (by norm_num : (20943712677888:ℕ) ≠ 0),
    Finsupp.add_apply, Nat.factorization_eq_zero_of_not_dvd (by decide : ¬ (5:ℕ) ∣ 20943712677888),
    Nat.factorization_pow, Finsupp.smul_apply,
    Nat.Prime.factorization (by norm_num : Nat.Prime 5), Finsupp.single_eq_same]; norm_num [nsmul_eq_mul]

/-- Non-square witness n = 15: `ν_3(15! * 15#) = 7` (odd); `15! * 15# = 3^7 * 17955857920000`, and `3 ∤ 17955857920000`. -/
theorem nonsquare_at_15 : ¬ IsSquare (15 ! * primorial 15) := by
  have hM : (15 : ℕ) ! * primorial 15 = 3 ^ 7 * 17955857920000 := by decide
  apply not_isSquare_of_odd_factorization (p := 3) (e := 3)
    (mul_ne_zero (Nat.factorial_ne_zero 15) (primorial_ne_zero 15))
  rw [hM, Nat.factorization_mul (by norm_num : ((3:ℕ)^7) ≠ 0) (by norm_num : (17955857920000:ℕ) ≠ 0),
    Finsupp.add_apply, Nat.factorization_eq_zero_of_not_dvd (by decide : ¬ (3:ℕ) ∣ 17955857920000),
    Nat.factorization_pow, Finsupp.smul_apply,
    Nat.Prime.factorization Nat.prime_three, Finsupp.single_eq_same]; norm_num [nsmul_eq_mul]

/-- Non-square witness n = 16: `ν_3(16! * 16#) = 7` (odd); `16! * 16# = 3^7 * 287293726720000`, and `3 ∤ 287293726720000`. -/
theorem nonsquare_at_16 : ¬ IsSquare (16 ! * primorial 16) := by
  have hM : (16 : ℕ) ! * primorial 16 = 3 ^ 7 * 287293726720000 := by decide
  apply not_isSquare_of_odd_factorization (p := 3) (e := 3)
    (mul_ne_zero (Nat.factorial_ne_zero 16) (primorial_ne_zero 16))
  rw [hM, Nat.factorization_mul (by norm_num : ((3:ℕ)^7) ≠ 0) (by norm_num : (287293726720000:ℕ) ≠ 0),
    Finsupp.add_apply, Nat.factorization_eq_zero_of_not_dvd (by decide : ¬ (3:ℕ) ∣ 287293726720000),
    Nat.factorization_pow, Finsupp.smul_apply,
    Nat.Prime.factorization Nat.prime_three, Finsupp.single_eq_same]; norm_num [nsmul_eq_mul]

/-- Non-square witness n = 17: `ν_3(17! * 17#) = 7` (odd); `17! * 17# = 3^7 * 83027887022080000`, and `3 ∤ 83027887022080000`. -/
theorem nonsquare_at_17 : ¬ IsSquare (17 ! * primorial 17) := by
  have hM : (17 : ℕ) ! * primorial 17 = 3 ^ 7 * 83027887022080000 := by decide
  apply not_isSquare_of_odd_factorization (p := 3) (e := 3)
    (mul_ne_zero (Nat.factorial_ne_zero 17) (primorial_ne_zero 17))
  rw [hM, Nat.factorization_mul (by norm_num : ((3:ℕ)^7) ≠ 0) (by norm_num : (83027887022080000:ℕ) ≠ 0),
    Finsupp.add_apply, Nat.factorization_eq_zero_of_not_dvd (by decide : ¬ (3:ℕ) ∣ 83027887022080000),
    Nat.factorization_pow, Finsupp.smul_apply,
    Nat.Prime.factorization Nat.prime_three, Finsupp.single_eq_same]; norm_num [nsmul_eq_mul]

/-- Non-square witness n = 18: `ν_2(18! * 18#) = 17` (odd); `18! * 18# = 2^17 * 24936491397943125`, and `2 ∤ 24936491397943125`. -/
theorem nonsquare_at_18 : ¬ IsSquare (18 ! * primorial 18) := by
  have hM : (18 : ℕ) ! * primorial 18 = 2 ^ 17 * 24936491397943125 := by decide
  apply not_isSquare_of_odd_factorization (p := 2) (e := 8)
    (mul_ne_zero (Nat.factorial_ne_zero 18) (primorial_ne_zero 18))
  rw [hM, Nat.factorization_mul (by norm_num : ((2:ℕ)^17) ≠ 0) (by norm_num : (24936491397943125:ℕ) ≠ 0),
    Finsupp.add_apply, Nat.factorization_eq_zero_of_not_dvd (by decide : ¬ (2:ℕ) ∣ 24936491397943125),
    Nat.factorization_pow, Finsupp.smul_apply,
    Nat.Prime.factorization Nat.prime_two, Finsupp.single_eq_same]; norm_num [nsmul_eq_mul]

/-- Non-square witness n = 19: `ν_2(19! * 19#) = 17` (odd); `19! * 19# = 2^17 * 9002073394657468125`, and `2 ∤ 9002073394657468125`. -/
theorem nonsquare_at_19 : ¬ IsSquare (19 ! * primorial 19) := by
  have hM : (19 : ℕ) ! * primorial 19 = 2 ^ 17 * 9002073394657468125 := by decide
  apply not_isSquare_of_odd_factorization (p := 2) (e := 8)
    (mul_ne_zero (Nat.factorial_ne_zero 19) (primorial_ne_zero 19))
  rw [hM, Nat.factorization_mul (by norm_num : ((2:ℕ)^17) ≠ 0) (by norm_num : (9002073394657468125:ℕ) ≠ 0),
    Finsupp.add_apply, Nat.factorization_eq_zero_of_not_dvd (by decide : ¬ (2:ℕ) ∣ 9002073394657468125),
    Nat.factorization_pow, Finsupp.smul_apply,
    Nat.Prime.factorization Nat.prime_two, Finsupp.single_eq_same]; norm_num [nsmul_eq_mul]

/-- Non-square witness n = 20: `ν_2(20! * 20#) = 19` (odd); `20! * 20# = 2^19 * 45010366973287340625`, and `2 ∤ 45010366973287340625`. -/
theorem nonsquare_at_20 : ¬ IsSquare (20 ! * primorial 20) := by
  have hM : (20 : ℕ) ! * primorial 20 = 2 ^ 19 * 45010366973287340625 := by decide
  apply not_isSquare_of_odd_factorization (p := 2) (e := 9)
    (mul_ne_zero (Nat.factorial_ne_zero 20) (primorial_ne_zero 20))
  rw [hM, Nat.factorization_mul (by norm_num : ((2:ℕ)^19) ≠ 0) (by norm_num : (45010366973287340625:ℕ) ≠ 0),
    Finsupp.add_apply, Nat.factorization_eq_zero_of_not_dvd (by decide : ¬ (2:ℕ) ∣ 45010366973287340625),
    Nat.factorization_pow, Finsupp.smul_apply,
    Nat.Prime.factorization Nat.prime_two, Finsupp.single_eq_same]; norm_num [nsmul_eq_mul]

/-- Non-square witness n = 21: `ν_2(21! * 21#) = 19` (odd); `21! * 21# = 2^19 * 945217706439034153125`, and `2 ∤ 945217706439034153125`. -/
theorem nonsquare_at_21 : ¬ IsSquare (21 ! * primorial 21) := by
  have hM : (21 : ℕ) ! * primorial 21 = 2 ^ 19 * 945217706439034153125 := by decide
  apply not_isSquare_of_odd_factorization (p := 2) (e := 9)
    (mul_ne_zero (Nat.factorial_ne_zero 21) (primorial_ne_zero 21))
  rw [hM, Nat.factorization_mul (by norm_num : ((2:ℕ)^19) ≠ 0) (by norm_num : (945217706439034153125:ℕ) ≠ 0),
    Finsupp.add_apply, Nat.factorization_eq_zero_of_not_dvd (by decide : ¬ (2:ℕ) ∣ 945217706439034153125),
    Nat.factorization_pow, Finsupp.smul_apply,
    Nat.Prime.factorization Nat.prime_two, Finsupp.single_eq_same]; norm_num [nsmul_eq_mul]

/-- Assembled form (cited by the ⟹ direction of the main theorem): `n! * n#` is non-square at every `n` with `6 ≤ n ≤ 21`. -/
theorem nonsquare_six_to_twenty_one (n : ℕ) (h6 : 6 ≤ n) (h21 : n ≤ 21) :
    ¬ IsSquare (n ! * primorial n) := by
  interval_cases n
  · exact nonsquare_at_6
  · exact nonsquare_at_7
  · exact nonsquare_at_8
  · exact nonsquare_at_9
  · exact nonsquare_at_10
  · exact nonsquare_at_11
  · exact nonsquare_at_12
  · exact nonsquare_at_13
  · exact nonsquare_at_14
  · exact nonsquare_at_15
  · exact nonsquare_at_16
  · exact nonsquare_at_17
  · exact nonsquare_at_18
  · exact nonsquare_at_19
  · exact nonsquare_at_20
  · exact nonsquare_at_21

/-- **Lemma B (interval prime lemma)**: for every integer `n ≥ 22` there exists a prime `p` with `n/3 < p ≤ n/2`.
    (`/` is natural-number division: the interval reads `⌊n/3⌋ < p ≤ ⌊n/2⌋`; for prime existence this coincides
    with the real reading, since `p` is an integer: `n/3 < p ⟺ ⌊n/3⌋ < p` and `p ≤ n/2 ⟺ p ≤ ⌊n/2⌋`.) -/
theorem exists_prime_between_div_three_and_div_two (n : ℕ) (hn : 22 ≤ n) :
    ∃ p : ℕ, p.Prime ∧ n / 3 < p ∧ p ≤ n / 2 := by
  rcases lt_or_ge n 1000000 with h | h
  · exact NFactPrimorialSquareChain.lemmaB_upto_999999 n hn (by omega)
  · exact NFactPrimorialSquareChain.lemmaB_tail_1000000 n h

/-- **Main theorem (complete solution set, iff form)**:
    `n! * n#` is a perfect square ⟺ `n ≤ 5` (for `n : ℕ`).
    Over ℕ, `n ≤ 5 ⟺ n ∈ {0,1,2,3,4,5}`, so this is the main statement "the solutions are exactly n ∈ {0,...,5}". -/
theorem isSquare_factorial_mul_primorial_iff (n : ℕ) :
    IsSquare (n ! * primorial n) ↔ n ≤ 5 := by
  constructor
  · -- ⟹ direction: n! * n# a square ⟹ n ≤ 5 (contrapositive: for n ≥ 6, n! * n# is not a square)
    intro h
    by_contra hle
    have h6 : 6 ≤ n := by omega
    rcases lt_or_ge n 22 with hlt | h22
    · -- the 6 ≤ n ≤ 21 segment: direct contradiction with the pointwise non-square witnesses of the witness layer
      exact nonsquare_six_to_twenty_one n h6 (by omega) h
    · -- the 22 ≤ n segment: Lemma B yields a prime p ∈ (n/3, n/2]; the criterion gives Odd (ν_p(n!)), but Legendre's sum evaluates ν_p(n!) = 2
      obtain ⟨p, hp, hpl, hpu⟩ := exists_prime_between_div_three_and_div_two n h22
      have hodd : Odd ((n !).factorization p) :=
        (isSquare_iff_all_factorization_odd n (by omega)).mp h p hp (by omega)
      have hn3p : n < 3 * p := by omega
      have hp3 : 3 ≤ p := by omega
      have hpsq : n < p ^ 2 := by
        have h3pp : 3 * p ≤ p * p := Nat.mul_le_mul hp3 (le_refl p)
        rw [pow_two]
        exact lt_of_lt_of_le hn3p h3pp
      have hlog2 : Nat.log p n < 2 := Nat.log_lt_of_lt_pow (by omega) hpsq
      have hd1 : n / p = 2 := by
        have hp0 : 0 < p := hp.pos
        have hr : n - 2 * p < p := by omega
        rw [show n = n - 2 * p + p * 2 by omega, Nat.add_mul_div_left _ _ hp0,
          Nat.div_eq_of_lt hr, zero_add]
      have hfact2 : (n !).factorization p = 2 := by
        have hIco : Finset.Ico 1 2 = ({1} : Finset ℕ) := by decide
        calc (n !).factorization p = ∑ i ∈ Finset.Ico 1 2, n / p ^ i :=
              Nat.factorization_factorial hp hlog2
          _ = ∑ i ∈ ({1} : Finset ℕ), n / p ^ i := by rw [hIco]
          _ = n / p ^ 1 := Finset.sum_singleton _ _
          _ = n / p := by rw [pow_one]
          _ = 2 := hd1
      rw [hfact2] at hodd
      obtain ⟨k, hk⟩ := hodd
      omega
  · -- ⟸ direction: for n ≤ 5, give the explicit square root pointwise (x = 1, 1, 2, 6, 12, 60)
    intro h
    obtain ⟨h0, h1, h2, h3, h4, h5⟩ := witnesses_zero_to_five
    interval_cases n
    · exact (isSquare_iff_exists_sq _).mpr ⟨1, h0.symm⟩
    · exact (isSquare_iff_exists_sq _).mpr ⟨1, h1.symm⟩
    · exact (isSquare_iff_exists_sq _).mpr ⟨2, h2.symm⟩
    · exact (isSquare_iff_exists_sq _).mpr ⟨6, h3.symm⟩
    · exact (isSquare_iff_exists_sq _).mpr ⟨12, h4.symm⟩
    · exact (isSquare_iff_exists_sq _).mpr ⟨60, h5.symm⟩

/-- **Main theorem (existence-witness form, delivered alongside the iff form)**:
    the Diophantine equation `x² = n! * n#` has a solution ⟺ `n ∈ {0,1,2,3,4,5}`. -/
theorem exists_sq_factorial_mul_primorial_iff (n : ℕ) :
    (∃ x : ℕ, x ^ 2 = n ! * primorial n) ↔ n ∈ ({0, 1, 2, 3, 4, 5} : Finset ℕ) := by
  have hbridge : (∃ x : ℕ, x ^ 2 = n ! * primorial n) ↔ IsSquare (n ! * primorial n) := by
    rw [isSquare_iff_exists_sq]
    constructor
    · rintro ⟨x, hx⟩
      exact ⟨x, hx.symm⟩
    · rintro ⟨r, hr⟩
      exact ⟨r, hr.symm⟩
  rw [hbridge, isSquare_factorial_mul_primorial_iff]
  simp only [Finset.mem_insert, Finset.mem_singleton]
  constructor <;> intro h <;> omega

end NFactPrimorialSquare
\end{lstlisting}

\section{Verification evidence}\label{sec:verif}

\begin{itemize}
\item Toolchain: Lean \texttt{v4.34.0} (\texttt{leanprover/lean4:v4.34.0}),
mathlib4 \texttt{v4.34.0}, revision \texttt{5ed29652}; the exact toolchain
pin ships as \texttt{proofs/lean-toolchain}.
\item Statement integrity: the frozen snapshot
\texttt{nfact-primorial-square.statement.txt} has sha256
\begin{center}
\texttt{0b13d08d6d915c2ebb25a070461a9d2a}\allowbreak
\texttt{7d0364b12b3627c1b1efba5cc9f8ea48},
\end{center}
the statement source \texttt{nfact-primorial-square.lean} has sha256
\begin{center}
\texttt{638064201e96a3bdb42e75b00020c16e}\allowbreak
\texttt{2fb27c2738c58fa3d4923cde2a451a8f},
\end{center}
and the final proof file \texttt{final-nfact-primorial-square.lean}
(1094 lines, LF) has sha256
\begin{center}
\texttt{276cb2e698f6d201752eed3cb1104254}\allowbreak
\texttt{d976f1a392ba3860a194eaa66986070a}.
\end{center}
The final audit compared all eight statement headers of the proof file
against the frozen snapshot under two whitespace conventions (eight of eight
match) and re-checked the kernel goal states against the snapshot's
\texttt{\#check} record (eight of eight match). The listing excerpts in this
note are byte-identical to the shipped proof files (mechanical record:
\texttt{audit/listing-verify.txt}). The three \texttt{proofs/} sources are
comment-cleaned editions (English docstrings): their code regions are
byte-identical to the frozen files listed above (correspondence:
\texttt{audit/comment-cleanup-diff.md}), and the post-cleanup kernel
re-verification is recorded in \texttt{audit/kernel-rerun.txt}.
\item Kernel check: strict compilation of the final file exits 0 with no
\texttt{sorry} and no \texttt{admit} (whole-file scan including comments:
zero hits; likewise zero hits for \texttt{axiom}, \texttt{unsafe},
\texttt{opaque}, \texttt{native\_decide}, \texttt{implemented\_by}). Three
diagnostics remain, all linter/deprecation-class and non-rejectable under
the audit rules applied: an elaborator exponentiation-threshold hint on
\texttt{18\^{}100000} (intentional; the kernel checks the proof term
independently), a deprecated \texttt{push\_neg} alias, and one
unused-variable linter note. \texttt{\#print axioms} output, verbatim,
re-printed independently by the packaging track on 2026-10-06 (strict compile
exit 0; the wrapped fourth entry is terminal formatting):
\begin{lstlisting}[basicstyle=\ttfamily\footnotesize,frame=none]
'NFactPrimorialSquare.isSquare_factorial_mul_primorial_iff' depends on axioms: [propext, Classical.choice, Quot.sound]
'NFactPrimorialSquare.exists_sq_factorial_mul_primorial_iff' depends on axioms: [propext, Classical.choice, Quot.sound]
'NFactPrimorialSquare.isSquare_iff_all_factorization_odd' depends on axioms: [propext, Classical.choice, Quot.sound]
'NFactPrimorialSquare.exists_prime_between_div_three_and_div_two' depends on axioms: [propext,
 Classical.choice,
 Quot.sound]
'NFactPrimorialSquare.witnesses_zero_to_five' depends on axioms: [propext, Classical.choice, Quot.sound]
'NFactPrimorialSquare.primorial_factorization_eq_one' depends on axioms: [propext, Classical.choice, Quot.sound]
'NFactPrimorialSquare.primorial_factorization_eq_zero' depends on axioms: [propext, Classical.choice, Quot.sound]
'NFactPrimorialSquare.factorization_factorial_mul_primorial' depends on axioms: [propext, Classical.choice, Quot.sound]
\end{lstlisting}
Every entry is exactly the whitelist \{propext, \texttt{Quot.sound},
\texttt{Classical.choice}\}; \texttt{sorryAx} does not occur.
\item Grading by the independent final audit (adjudication of 2026-10-06,
four checks all passed): \emph{complete proof}. A cleaned copy of the
adjudication ships as \texttt{audit/gate3-final-adjudication.txt}
(translated to English; machine-local path prefixes replaced by
\texttt{<REPO-ROOT>}; verdict text and data unchanged --- normalization
disclosed in \texttt{audit/package-notes.md}); the well-definedness gate
output ships verbatim as \texttt{audit/welldef-nfact-primorial-square.txt}.
\item Audit chain: statement frozen by hash before proving; post-hoc
symmetric-difference comparison (text-level and kernel goal-state level);
final strict recompile; axiom re-print (this note's copy:
\texttt{audit/axioms.out}); novelty record \texttt{audit/c9-record.md}. The
claim audit \texttt{claim-check.md} is added to \texttt{audit/} after the
independent claim audit returns its adjudication and is a prerequisite for
publication.
\end{itemize}

\section{Novelty evidence}\label{sec:novelty}

Adjudication carried over from the independent novelty review:
\textbf{no occupying record found}; search-strength label \textbf{strong}.
The closing dates of the search are 2026-10-05 (deep-dig layer) and
2026-10-06 (independent review layer, which maintained every load-bearing
claim). The full query log, including every zero-hit query, ships in full as
\texttt{audit/c9-record.md} (English cleaned edition; conclusions and query
data unedited). A zero hit is not novelty; the strongest claim licensed here
is that no record of this object was found in the channels queried.

The problem is on record as open in the source paper. Novakovi\'c writes
(verbatim wording; mathematical notation retypeset): ``We believe that
finitely many solutions occur for $x^2 = n! \cdot n\#$ or
$x^2 = n!! \cdot [1, \ldots, n]$, but at the moment we were not able to give
an argument'' \cite{novakovic2026}. The present note answers the first of
these two cases completely; the second case is \emph{not} claimed here.

\begin{center}\small
\begin{tabular}{p{3.4cm}p{5.4cm}p{6.0cm}}
\hline
channel & queries & result \\
\hline
literature (author self-report) & arXiv:2601.16757 body text & the case is on
record as open (quotation above) \cite{novakovic2026} \\
arXiv metadata & primorial $\wedge$ factorial $\wedge$ perfect square;
primorial $\wedge$ factorial $\wedge$ square; double factorial $\wedge$
least common multiple; squarefree part $\wedge$ factorial & zero hits each \\
arXiv metadata & Brocard $\wedge$ primorial & one hit \cite{peixoto2026}:
adjacent object (primorial-core Pell), does not characterize this equation \\
OpenAlex full text & brocard/primorial/square/factorial family; k-free
Brocard family & nearest hits include
\cite{novakovic2026,novakovic2026b}; no same-form statement \\
OEIS keyword & ``n! primorial square''; ``factorial times primorial is a
square'' & hits exist; no same-form \\
OEIS keyword & ``double factorial lcm square''; ``Novakovic factorial'' &
zero hits \\
public Lean ecosystem (formal-conjectures tarball, content level) &
\texttt{primorial}; \texttt{doubleFactorial} & occurrences in unrelated
files only; no item in this domain \\
OEIS value columns (review layer) & $x$: 1,1,2,6,12,60; $x^2$: 1,4,36,144,3600
& no registration of these objects \\
local mathlib (pinned v4.34.0), positive control & content scan with live
control & channel live; no same-form item \\
\hline
\end{tabular}
\end{center}

Delimitations against adjacent records (wording obligations carried from the
source task card):
(i) OEIS A055204 \cite{oeisA055204} asks for the \emph{minimal} $a$ making
$n! \cdot a$ a square; the present theorem asks the \emph{specific} value
$a = n\#$ and is not an instance of A055204.
(ii) Erd\H{o}s problem \#399 \cite{erdos399} concerns $n! = x^k \pm y^k$ --- a
different equation from $x^2 = n! \cdot n\#$ --- and is marked DISPROVED
(2026-02-04, counterexample $10! = 48^4 - 36^4$); it is recorded here only as
same-family adjacency, never as an open anchor.
(iii) arXiv:2602.09687 \cite{novakovic2026b} is recorded with the restricted
wording that \emph{its object of study does not contain} $n\#$.
(iv) arXiv:2606.23485 \cite{peixoto2026} (\S 5.3.1, primorial-core Pell
equation $X^2 - D Z^2 = 1$) is an adjacent object and does not characterize
which $n$ make $n! \cdot n\#$ a square.
(v) arXiv:2610.01899 \cite{yudin2026} (square products of factorials and a
conjecture of Erd\H{o}s--Graham) is a different object; no collision, but a
citation duty is recorded and discharged here.
(vi) OEIS A166968 \cite{oeisA166968} (whose $a(2) = 8$ says: for $m \ge 8$
there is a prime in $(m, 3m/2]$) provides a literature-level shortest path
toward Lemma~\ref{lem:B} in the $m$-parametrization; the kernel proof shipped
here does not go through it.
The failure set of the interval property below $n = 22$ is $\{2,3,9,21\}$
under the domain convention $n \ge 2$ ($n = 0, 1$ also fail, on an empty
interval; this is irrelevant to Lemma~\ref{lem:B}, which starts at
$n = 22$).

\section{Scope and limitations}\label{sec:limits}

Grade and boundary, stated as the audit states them. (a) The kernel layer
proves exactly the eight frozen statements of Listing~\ref{lst:stmt}; the
grade of each is \emph{complete proof} (zero \texttt{sorry}, axioms within
the whitelist, statement text unchanged from the frozen snapshot, strict
compile exit 0). Nothing beyond is claimed. (b) The second equation of the
same remark, $x^2 = n!! \cdot [1, \ldots, n]$, is \emph{not} claimed: only an
exclusion path for even $n$ was obtained. (c) No OEIS
registration is claimed or performed by this deposit. (d) The exclusion of
$n \ge 22$ is a pure theorem (Lemma~\ref{lem:B}) and does not rely on the
Berndt--Galway computational bound ($10^9$) on the Brocard side; that bound
plays no role here. (e) The hypothesis $n \ge 1$ of Lemma~\ref{lem:A} is not
used by its proof (the criterion holds at $n = 0$ as well); the lemma is
stated as frozen. (f) A list of background items recorded as undecided at
the full-text level is enumerated in the source task card (which stays in the
source repository and is not shipped); none of them is load-bearing for this
result and none is cited here as a zero hit. (g) This note contains no
literature review; it is a deposit record, not a survey paper.

\section*{Data availability}

Lean sources, verification and novelty-search evidence are deposited on
Zenodo (concept DOI \url{https://doi.org/10.5281/zenodo.23177003}; the
version history is recorded in \texttt{metadata/README.md} of the deposit),
with a public mirror of the deposit contents on GitHub at
\url{https://github.com/Aurora0134/ai4math-claims}. The deposit contains the
claim note (LaTeX source and PDF), the Lean sources (\texttt{proofs/}:
statement snapshot, statement source, final proof file,
\texttt{lean-toolchain}), and the audit and novelty-search records
(\texttt{audit/}). Claim note under CC BY 4.0, code under MIT.

\section*{Use of generative AI}

The formalization and proofs reported in this note were developed with AI
assistance (large language model: \texttt{kimi-k3}) under human direction
and review. All statements and proofs reported here were checked by the Lean
4 kernel (Lean v4.34.0, mathlib v4.34.0): every proof compiles with no
\texttt{sorry}, and \texttt{\#print axioms} reports only subsets of
\{propext, Quot.sound, Classical.choice\}. The author reviewed all content
and takes full responsibility for it.

\section*{Author contributions}

The author selected and adjudicated the topic, verified the semantics of the
formal statements, signed off the verification gates, and decided to deposit.

\begin{thebibliography}{9}
% Only works actually cited; each was re-checked to exist before entry
% (2026-10-06: lean4/mathlib via live Crossref title+pages metadata;
% the four arXiv items via the live abs-page citation_title/citation_author/
% citation_date metadata; erdos399 via the live page, whose status widget
% reads DISPROVED; the two OEIS entries via the live OEIS API name fields).
% No entry is quoted from memory. Record: audit/bib-verify.md (ships).
\bibitem{lean4}
L.~de Moura and S.~Ullrich.
\newblock The Lean 4 theorem prover and programming language.
\newblock \emph{Automated Deduction -- CADE-28}, LNCS, pp.~625--635, 2021.
\url{https://doi.org/10.1007/978-3-030-79876-5_37}

\bibitem{mathlib}
The mathlib Community.
\newblock The Lean mathematical library.
\newblock \emph{CPP 2020}, pp.~367--381, 2020.
\url{https://doi.org/10.1145/3372885.3373824}

\bibitem{novakovic2026}
S.~Novakovi\'c.
\newblock The Diophantine equation $P(x) = \prod_{i=1}^{r} H_{n_i}$.
\newblock arXiv:2601.16757, 2026. \url{https://arxiv.org/abs/2601.16757}

\bibitem{novakovic2026b}
S.~Novakovi\'c.
\newblock On a generalization of the Brocard--Ramanujan Diophantine equation.
\newblock arXiv:2602.09687, 2026. \url{https://arxiv.org/abs/2602.09687}

\bibitem{peixoto2026}
L.~V. Peixoto.
\newblock The Brocard problem: structural invariants, p-adic density, and the
generative sieve.
\newblock arXiv:2606.23485, 2026. \url{https://arxiv.org/abs/2606.23485}

\bibitem{yudin2026}
F.~Yudin.
\newblock Square products of factorials and a conjecture of Erd\H{o}s and
Graham.
\newblock arXiv:2610.01899, 2026. \url{https://arxiv.org/abs/2610.01899}

\bibitem{erdos399}
Erd\H{o}s Problems website.
\newblock Problem \#399 ($n! = x^k \pm y^k$; marked DISPROVED, counterexample
$10! = 48^4 - 36^4$).
\newblock \url{https://www.erdosproblems.com/399}

\bibitem{oeisA055204}
OEIS Foundation.
\newblock Squarefree part of $n!$: $n!$ divided by its largest square divisor
(sequence A055204).
\newblock \emph{The On-Line Encyclopedia of Integer Sequences}.
\url{https://oeis.org/A055204}

\bibitem{oeisA166968}
OEIS Foundation.
\newblock Minimum $k$ such that for all $m \ge k$ there is a prime $p$ with
$m < p < m(n+1)/n$ (sequence A166968).
\newblock \emph{The On-Line Encyclopedia of Integer Sequences}.
\url{https://oeis.org/A166968}
\end{thebibliography}

\end{document}
